\section{A fresh quadratic host preserving a finite scalar witness}
\label{sec:cubic-fresh-quadratic-host}

\begin{proposition}[AC-G77: an unramified CM host with the required split places]
\label{prop:ACG77}
Let \(K=\mathbf Q(\sqrt{-3})\), let \(p>3\), and let \(S\) be a finite set of rational primes containing \(2,3,p,269\). Let \(\Omega/\mathbf Q\) be finite Galois, contain \(K\), and be unramified outside \(S\). Any finite residual data and finite characters to be preserved may be included in \(\Omega\).

There are arbitrarily large primes \(r\notin S\) such that
\[
 K'=\mathbf Q(\sqrt{-r}),\qquad F^+=\mathbf Q(\sqrt{3r}),\qquad
 F=KK'=F^+K'
 \tag{G77.1}
\]
have the following properties.
\begin{enumerate}
 \item \(F/F^+\) is CM quadratic and unramified at every finite place;
 \([F:\mathbf Q]=4\), \([F^+:\mathbf Q]=2\), and
 \[
 \operatorname{disc}(F^+)=3r,\qquad
 |\operatorname{disc}(F)|=(3r)^2.
 \tag{G77.2}
 \]
 \item Every \(q\in S\) splits in \(K'\), and every place of \(F^+\) above \(q\) splits in \(F/F^+\). Thus \(F_0=K'\) has \(p\) split even if \(p\) is inert in \(K\).
 \item \(K'\cap\Omega=\mathbf Q\), \(F\cap\Omega=K\), and every
 \(\sigma\in\operatorname{Gal}(\Omega/K)\) has an extension fixing \(F\) and preserving exactly its action on \(\Omega\).
 If \(\sigma\) has order prime to \(p\), this extension can fix the entire cyclotomic \(\mathbf Z_p\)-extension \(F_\infty/F\).
 \item \(F\) contains no primitive \(p\)-th root of unity, and \(\chi_{269}|_{G_F}\) still has order three.
\end{enumerate}
An \emph{existing} witness, for example
\(\bar V(\sigma)=-I_2\), \(\bar\kappa(\sigma)=1\), together with specified finite-character equations, is therefore retained. Its existence is a premise.
\end{proposition}

\begin{proof}
Put \(m=8\prod_{q\in S,\ q\ {\rm odd}}q\). Dirichlet's theorem gives arbitrarily large primes \(r\equiv-1\pmod m\). They are outside \(S\), and \(r\equiv7\pmod8\), so \(K'\) has fundamental discriminant \(-r\). At every odd \(q\in S\), \(-r\) is a square in \(\mathbf Q_q\), being \(1\) modulo \(q\). It is a square in \(\mathbf Q_2\) as well, being \(1\) modulo \(8\). Thus all primes of \(S\) split in \(K'\).

The two distinct imaginary quadratic fields have biquadratic compositum \(F\). Complex conjugation changes both square roots and fixes exactly \(F^+\). The squarefree integer \(3r\) is \(1\) modulo \(4\), giving \(\operatorname{disc}(F^+)=3r\).

Outside \(3,r\), both quadratic fields are unramified, so \(F/F^+\) is unramified. At \(3\), \(K'\) splits and \(F/F^+\) splits. At \(r\), \(K/\mathbf Q\) is unramified, and adjoining this unramified quadratic local algebra to \(F^+\) is unramified. Thus the relative discriminant is the unit ideal, and discriminant transitivity gives (G77.2). More generally, for \(q\in S\),
\[
 F\otimes\mathbf Q_q=(K\otimes\mathbf Q_q)\times(K\otimes\mathbf Q_q),
 \qquad F^+\otimes\mathbf Q_q=K\otimes\mathbf Q_q.
 \tag{G77.3}
\]
This proves the required split-place assertions, including when \(K\) has a quadratic \(p\)-adic completion.

The field \(K'\) ramifies at \(r\), whereas \(\Omega\) does not. Therefore \(K'\cap\Omega=\mathbf Q\). The product description of
\(\operatorname{Gal}(\Omega K'/\mathbf Q)\) extends \(\sigma\) by the identity on \(K'\). It fixes \(K,K'\), hence \(F\), and also shows \(F\cap\Omega=K\).

If \(\sigma\) has order prime to \(p\), so does its extension. The intersection
\(\Omega K'\cap F_\infty\) is a finite Galois \(p\)-extension of \(F\), so that automorphism acts trivially on the intersection. Combine it with the identity on \(F_\infty\), then lift to an algebraic closure. Its finite equations on \(\Omega\) remain unchanged.

For scalar and finite-character equations whose specified values have order prime to \(p\), one may first replace a witness by its prime-to-\(p\) part. Explicitly, if its order is \(p^a b\) with \((p,b)=1\), choose \(t\equiv0\pmod{p^a}\), \(t\equiv1\pmod b\). The element \(\sigma^t\) has order prime to \(p\) and preserves all those specified values. This applies to \(-I_2\), \(\bar\kappa=1\), and the finite characters here. It preserves these equations, not the original automorphism on every element of \(\Omega\).

Every quotient of \(\operatorname{Gal}(F/\mathbf Q)\) has exponent at most two. An inclusion \(\mu_p\subset F\) would give a cyclic quotient of order \(p-1\), impossible for \(p>3\). Finally an order-three character retains its image on a subgroup of index two. This proves the last assertion.
\end{proof}

\begin{corollary}[AC-G77-restriction: irreducibility and the packet field hypotheses]
\label{cor:ACG77restriction}
Let \(W\) be an absolutely irreducible two-dimensional representation of \(G_K\), in characteristic different from two, unramified at all places above the chosen \(r\). Then \(W|_{G_F}\) is absolutely irreducible.

Suppose an automorphic block datum is unramified outside \(S\), its finite local parameters remain unramified after base change at unramified places, and its pair of normalized archimedean weights is parallel. On (G77.1), every ramified finite place splits over \(F^+\), \(p\) splits in \(F_0=K'\), and each weight pair above either embedding of \(F_0\) occurs twice. These supply the field, split-place and even-multiplicity requirements in Thorne's Theorem 4.5. The automorphic blocks, conjugate self-duality, regular algebraicity, and ordinary support remain separate inputs.
\end{corollary}

\begin{proof}
Let \(\epsilon\) be the character of \(G_K/G_F\). It ramifies at every place above \(r\). If \(W|_{G_F}\) had an invariant line, its conjugate under the other coset would be a distinct line, since \(W\) is irreducible over \(G_K\). The two lines form a direct sum. The operator \(1,-1\) on them commutes with \(G_F\) and anticommutes with the other coset, giving \(W\simeq W\otimes\epsilon\). At inertia above \(r\) this would identify the identity with \(-I_2\), impossible in characteristic different from two. The argument also works over an algebraic closure of coefficients.

An unramified finite local parameter stays unramified under restriction. Thus the stipulated block datum acquires no ramification outside \(S\), including at \(r\). Formula (G77.3) treats all its ramified places. Exactly two embeddings of \(F\) extend a fixed embedding of \(K'\); parallel normalized weight pairs are equal at these embeddings and hence have even multiplicity. This proves exactly the asserted field interfaces.
\end{proof}

\paragraph{Sources, dictionary, and scope.}
Dirichlet's theorem is used only for existence of \(r\); its statement appears in Milne, \emph{Class Field Theory}, version 4.03, Introduction, \emph{Density theorems}, printed p.\ 7. No effective upper bound for \(r\) is asserted. Thorne, \emph{A \(p\)-adic approach to the existence of level-raising congruences}, \emph{Proc.\ London Math.\ Soc.} 128 (2024), e12584, \S4, Theorem 4.5, printed pp.\ 21--23, was read with its complete printed proof; the deeper cited trace-formula inputs are retained. Its parity condition concerns the weights of \(\pi_i|\cdot|^{(n_i-n)/2}\), not unnormalized block weights.
The finite unramified map \(\operatorname{Spec}\mathcal O_F\to
\operatorname{Spec}\mathcal O_{F^+}\) and coefficient restrictions fit the
already-audited SGA~1, Expose~V, \S\S4 and 7 framework, with no new archive
reading. This construction preserves an existing witness and supplies
field data; it does not create the witness, a stable three-dimensional
packet, or the actual characteristic-\(p\) support containment.
